\documentclass[12pt,a4paper]{article}

\usepackage[utf8]{inputenc}
\usepackage[T1]{fontenc}
\usepackage{lmodern}
\usepackage{amsmath,amssymb,amsthm}
\usepackage[top=2.5cm,bottom=2.5cm,left=2.8cm,right=2.8cm]{geometry}
\usepackage{parskip}
\usepackage{booktabs}

\newtheorem{theorem}{Théorème}
\newtheorem{lemma}[theorem]{Lemme}
\newtheorem{corollary}[theorem]{Corollaire}
\newtheorem{remark}[theorem]{Remarque}
\newtheorem{definition}[theorem]{Définition}
\newtheorem{example}[theorem]{Exemple}

\newcommand{\Inc}{\ensuremath{\mathrm{Inc}}}
\newcommand{\Exc}{\ensuremath{\mathrm{Exc}}}
\newcommand{\vbar}{\bar{v}}

\begin{document}

\section*{Analyse théorique de la règle à 4 cas}

Cette section prouve la convergence de la règle locale à 4 cas elle-même et pas le systeme complet : 

\subsection*{}

Soit $(X_t,Y_t)_{t\ge1}\overset{\text{i.i.d.}}{\sim}P$ de loi jointe $P(X,Y)$
arbitraire. Deux automates à $2N$ états gouvernent le littéral $x$ et sa
négation : $v$ et $\vbar$, avec $\mathrm{included}(s)=\mathbf 1[s>N]$.

\begin{center}
\begin{tabular}{cccc}
\toprule
$x$ & $y$ & Action sur $v$ & Action sur $\vbar$ \\
\midrule
0 & 0 & vers $\Inc(S)$ & vers $\Exc(S)$ \\
0 & 1 & vers $\Exc(a)$ & -- \\
1 & 0 & vers $\Exc(S)$ & vers $\Inc(S)$ \\
1 & 1 & -- & vers $\Exc(a)$ \\
\bottomrule
\end{tabular}
\end{center}

avec $S,a\in(0,1]$. Les probabilités de transition intérieures s'écrivent, pour toute loi $P$ :
\begin{align}
  p^+ &= S\,P(0,0), & p^- &= S\,P(1,0) + a\,P(0,1); \label{eq:pv}\\
  \bar p^+ &= S\,P(1,0), & \bar p^- &= S\,P(0,0) + a\,P(1,1). \label{eq:pvbar}
\end{align}



\subsection*{Loi stationnaire exacte}

\begin{theorem}\label{thm:stat}
Soit $\rho=p^+/p^-$, avec $p^+,p^->0$. La chaîne $(v_t)$ est irréductible et
apériodique sur $\{1,\ldots,2N\}$, donc $\mathcal L(v_t)\xrightarrow[t\to\infty]{}\pi$
(convergence en loi vers l'unique loi stationnaire, théorème ergodique
standard pour les chaînes de Markov finies irréductibles apériodiques), de
loi stationnaire
\[
  \pi(k) = \frac{(1-\rho)\rho^{k-1}}{1-\rho^{2N}}\ (\rho\ne1), \qquad
  \pi(k) = \frac1{2N}\ (\rho=1), \qquad
  \pi(\Inc) = \frac{\rho^N}{1+\rho^N}.
\]
Si $\rho>1$ : $\pi(\Inc)\to1$ et $1-\pi(\Inc)\le\rho^{-N}$. Si $\rho<1$ :
$\pi(\Exc)\to1$ symétriquement. De plus $\frac1T\sum_{t\le T}\mathbf1[v_t>N]\to\pi(\Inc)$ p.s.
\emph{Attention} : ceci est une convergence en loi (stationnaire) et une
convergence p.s.\ de la \emph{fréquence temporelle} -- pas une convergence
p.s.\ de $v_t$ vers un état fixe. Dans le cas ergodique ($p^+,p^->0$),
l'automate continue de se déplacer indéfiniment ; il n'y a stabilisation
permanente que dans le cas absorbant (Corollaire~\ref{cor:absorb}).
\end{theorem}

\begin{proof}
\emph{Irréductibilité.} Chaîne de naissance-mort à transitions intérieures
strictement positives.

\emph{Noyau explicite.} Pour $2\le k\le 2N-1$ : $P(k,k{+}1)=p^+$,
$P(k,k{-}1)=p^-$, $P(k,k)=1-p^+-p^-$. Aux bornes (troncature) :
$P(1,1)=1-p^+$, $P(1,2)=p^+$, $P(2N,2N)=1-p^-$, $P(2N,2N{-}1)=p^-$.

\emph{Apériodicité.} $P(1,1)=1-p^+\ge p^->0$ : l'état $1$ a un cycle de
longueur $1$, donc une période $d(1)=1$. Par irréductibilité, la période est
une propriété de classe : toute la chaîne est apériodique.

\emph{Forme géométrique.} Bilan détaillé $\pi(k)p^+=\pi(k{+}1)p^-$ donne
$\pi(k{+}1)/\pi(k)=\rho$ constant, soit $\pi(k)=\pi(1)\rho^{k-1}$ ;
normalisation $\sum_{k=1}^{2N}\pi(k)=1=\pi(1)\sum_{k=0}^{2N-1}\rho^k$ donne
$\pi(1)=(1-\rho)/(1-\rho^{2N})$ (série géométrique finie), d'où les formules
annoncées. $\pi(\Inc)=\sum_{k=N+1}^{2N}\pi(k)=\pi(1)\rho^N\sum_{k=0}^{N-1}\rho^k
=\pi(1)\rho^N(1-\rho^N)/(1-\rho)=\rho^N(1-\rho^N)/(1-\rho^{2N})=\rho^N/(1+\rho^N)$.
La convergence p.s.\ de la moyenne temporelle est le théorème ergodique
standard.
\end{proof}

\begin{corollary}[Cas absorbant]\label{cor:absorb}
Si $p^-=0,p^+>0$ : absorption p.s. en $2N$ en temps fini,
$\mathbb E[T_{2N}]=(2N-s_0)/p^+$ ; première entrée dans $\Inc$ (franchissement
de $N$), généralement plus rapide, $\mathbb E[T_{\Inc}]=(N+1-s_0)_+/p^+$. Si
$p^+=0,p^->0$ : symétrique en $1$/$\Exc$. Si $p^+=p^-=0$ : aucune transition
ne se produit jamais (chaque état trivialement absorbant), l'automate reste
figé à $s_0$.
\end{corollary}

\subsection*{Exclusion des dérives contradictoires}

\begin{lemma}\label{lem:excl}
Soit $\Delta=p^+-p^-$, $\bar\Delta=\bar p^+-\bar p^-$. Pour toute loi $P(X,Y)$ :
\[
  \Delta+\bar\Delta = -a\,P(Y{=}1) \le 0.
\]
Si $P(Y{=}1)>0$, l'inégalité est \emph{stricte} : $\Delta+\bar\Delta<0$, donc
$\Delta=\bar\Delta=0$ simultanément est impossible. Si $P(Y{=}1)=0$ (cas
dégénéré), $\Delta+\bar\Delta=0$ exactement.
\end{lemma}

\begin{proof}
De \eqref{eq:pv}--\eqref{eq:pvbar} : $\Delta = S\,P(0,0)-S\,P(1,0)-a\,P(0,1)$,
$\bar\Delta = S\,P(1,0)-S\,P(0,0)-a\,P(1,1)$. La somme élimine les termes en
$S$ : $\Delta+\bar\Delta=-a[P(0,1)+P(1,1)]=-a\,P(Y{=}1)\le0$, avec égalité
ssi $P(Y{=}1)=0$ (puisque $a>0$).
\end{proof}

\begin{corollary}\label{cor:asym}
Si $\Delta>0$ alors $\bar\Delta=-a\,P(Y{=}1)-\Delta<0$ strictement.
\end{corollary}

\begin{remark}
Le lemme exclut $\Delta>0$ et $\bar\Delta>0$ simultanément ; il n'exclut pas
$\Delta<0$ et $\bar\Delta<0$ simultanément (double exclusion, régime neutre).
Pour $N$ fini et $\rho\ne1$, l'état instantané peut néanmoins se trouver du
mauvais côté avec probabilité $\le\rho^{-N}\to0$.
\end{remark}

\subsection*{Sélection sous loi arbitraire}

\begin{theorem}\label{thm:sel}
Supposons $\Delta\ne0$ et $\bar\Delta\ne0$. Notons $\gamma=\max(\rho,1/\rho)>1$,
$\bar\gamma=\max(\bar\rho,1/\bar\rho)>1$ (probabilité du mauvais côté toujours
$\le\gamma^{-N}$, resp.\ $\bar\gamma^{-N}$, quel que soit le signe de
$\Delta,\bar\Delta$). Alors exactement un des trois régimes suivants a lieu :
\begin{itemize}
  \item $\Delta>0$ (donc $\bar\Delta<0$) $\iff S\,P(0,0)>S\,P(1,0)+a\,P(0,1)$ :
  \[
    \lim_{N\to\infty}\;\limsup_{t\to\infty}\;\mathbb P\bigl((v_t,\vbar_t)\ne(\Inc,\Exc)\bigr) = 0.
  \]
  \item $\bar\Delta>0$ (donc $\Delta<0$) $\iff S\,P(1,0)>S\,P(0,0)+a\,P(1,1)$ :
  \[
    \lim_{N\to\infty}\;\limsup_{t\to\infty}\;\mathbb P\bigl((v_t,\vbar_t)\ne(\Exc,\Inc)\bigr) = 0.
  \]
  \item $\Delta<0$ et $\bar\Delta<0$ :
  \[
    \lim_{N\to\infty}\;\limsup_{t\to\infty}\;\mathbb P\bigl((v_t,\vbar_t)\ne(\Exc,\Exc)\bigr) = 0.
  \]
\end{itemize}

\end{theorem}

\begin{proof}
Les équivalences découlent de \eqref{eq:pv}--\eqref{eq:pvbar}. L'exclusion
mutuelle des trois régimes découle du Corollaire~\ref{cor:asym}. Passage du
marginal au joint (cas $\Delta>0$ : $\rho>1$ donc $\gamma=\rho$, $\bar\rho<1$
donc $\bar\gamma=1/\bar\rho$ ; similaire ailleurs), \emph{sans} invoquer de loi
stationnaire jointe (la chaîne $(v_t,\vbar_t)$ n'est pas nécessairement
irréductible -- ex.\ $S{=}1$, $P(Y{=}1){=}0$ conserve $v_t+\vbar_t$). En
prenant d'abord $\limsup_{t\to\infty}$ (la borne du Théorème~\ref{thm:stat}
porte sur la loi \emph{stationnaire}, atteinte seulement asymptotiquement en
$t$ -- fausse pour un $t$ petit et une initialisation arbitraire), puis par
union bound,
\[
  \limsup_{t\to\infty}\mathbb P\bigl((v_t,\vbar_t)\ne(\Inc,\Exc)\bigr)
  \le \limsup_{t\to\infty}\mathbb P(v_t\ne\Inc)+\limsup_{t\to\infty}\mathbb P(\vbar_t\ne\Exc)
  \le \gamma^{-N}+\bar\gamma^{-N},
\]
sans hypothèse d'indépendance entre $v_t$ et $\vbar_t$. Puis
$\gamma^{-N}+\bar\gamma^{-N}\xrightarrow[N\to\infty]{}0$ donne la double
limite annoncée. Noter $\bar\gamma^{-N}=\bar\rho^N$ ici (pas $\bar\rho^{-N}$,
puisque $\bar\rho<1$) : c'est pour éviter cette confusion de signe que
$\gamma=\max(\rho,1/\rho)$ est défini une fois pour toutes.
\end{proof}



\subsection*{Concentration multi-bit}

\begin{theorem}\label{thm:multi}
Clause à $n$ features, $K$ fixées par la condition, $d=n-K$ libres gouvernées par
$2d$ automates. Pour chaque automate $j$ (préférence stationnaire non
dégénérée), $r_j=p_j^+/p_j^-$, $\gamma_j=\max(r_j,1/r_j)>1$,
$\gamma_{\min}=\min_j\gamma_j$. Alors, sans hypothèse d'indépendance,
l'erreur de \emph{récupération structurelle} vérifie :
\[
  \limsup_{t\to\infty}\mathbb P(C_t\ne C^\star) \le \sum_{j=1}^{2d}\frac1{1+\gamma_j^N}
  \le 2(n-K)\,\gamma_{\min}^{-N} \xrightarrow[N\to\infty]{} 0.
\]
\end{theorem}

\begin{proof}
Soit $A_j$ l'évènement « l'automate $j$ est du mauvais côté ». Si tous les
$A_j$ sont faux, chaque automate est dans son état préféré, donc $C_t=C^\star$ ;
par contraposée, $\{C_t\ne C^\star\}\subseteq A_1\cup\cdots\cup A_{2d}$, d'où
par l'union bound (Boole) :
\[
  \mathbb P(C_t\ne C^\star)\le\sum_{j=1}^{2d}\mathbb P(A_j).
\]
Par le Théorème~\ref{thm:stat}, $\mathbb P(A_j)=\pi(\Exc)=1/(1+r_j^N)$ si
$r_j>1$, ou $\pi(\Inc)=r_j^N/(1+r_j^N)$ si $r_j<1$ -- dans les deux cas
exactement $1/(1+\gamma_j^N)$, d'où la première inégalité. Ensuite,
$1/(1+\gamma_j^N)<\gamma_j^{-N}$ (dénominateur plus grand), et
$\gamma_j\ge\gamma_{\min}$ donne $\gamma_j^{-N}\le\gamma_{\min}^{-N}$ ; en
remplaçant chacun des $2d$ termes par cette même borne on obtient
$2(n-K)\,\gamma_{\min}^{-N}$.
\end{proof}





\section*{Application : ensemble conditionné et boosté}

On intègre la règle à 4 cas dans un ensemble conditionné (gates aléatoires)
et boosté (AdaBoost) pour illustrer son utilité pratique en classification.
\textbf{L'analyse de convergence complète de cet ensemble --
couverture des conditions sous la distribution adaptative $Q_m$, et
convergence du vote pondéré final -- n'est pas prouvée ici et reste un
travail futur.} Les résultats ci-dessous (couverture, boosting) caractérisent
des propriétés partielles utiles mais ne ferment pas cette boucle.

\subsection*{Algorithme}

Classifieur multiclasse one-vs-rest : une classe $c$ est traitée par
l'algorithme ci-dessous, indépendamment des autres.

\begin{quote}\ttfamily\small
\textbf{Entrée} : dataset $\{(x_i,y_i)\}$, $n$ features, $N$ états, $M$
clauses, $K$ features gelées/clause, $total$ tirages/clause, $S,a$.\\[4pt]
$w_i \gets 1/|\text{train}|$ pour tout $i$\\[2pt]
\textbf{pour} $m = 1$ à $M$ :\\
\hspace*{1.5em}tirer $F\subseteq\{1,\ldots,n\}$, $|F|=K$, uniformément (shuffle)\\
\hspace*{1.5em}tirer $v\in\{0,1\}^K$ uniformément (pile-ou-face indépendant/feature)\\
\hspace*{1.5em}créer une clause $C_m$ conditionnée sur $(F,v)$ ; geler ces $K$ littéraux\\
\hspace*{1.5em}$S_m \gets \{i : (x_i)_F = v\}$\\
\hspace*{1.5em}\textbf{si} $S_m=\emptyset$ : $\alpha_m \gets 0$ ; \textbf{continuer}\\
\hspace*{1.5em}\textbf{pour} $t = 1$ à $total$ :\\
\hspace*{3em}tirer $i\in S_m$ selon $Q_m$ (50/50 classe, puis $\propto w_i$ dans la classe)\\
\hspace*{3em}mettre à jour $C_m$ avec $(x_i,\mathbf 1[y_i{=}c])$ (règle à 4 cas)\\
\hspace*{1.5em}\textbf{si} $C_m$ vide (aucun littéral inclus) : $\alpha_m \gets 0$ ; \textbf{continuer}
\hspace*{1em}(pas de mise à jour des poids)\\
\hspace*{1.5em}$\mathrm{err}_m \gets \sum_{i\in S_m} w_i\mathbf 1[C_m(x_i)\ne\mathbf 1[y_i{=}c]] \big/ \sum_{i\in S_m} w_i$ \hspace*{0.5em}(sous $D_m$)\\
\hspace*{1.5em}$\mathrm{err}_m \gets \mathrm{clip}(\mathrm{err}_m,\epsilon,1-\epsilon)$\\
\hspace*{1.5em}$\alpha_m \gets \tfrac12\ln\bigl((1-\mathrm{err}_m)/\mathrm{err}_m\bigr)$\\
\hspace*{1.5em}\textbf{pour} $i \in S_m$ : $w_i \gets w_i \cdot \exp(\pm\alpha_m)$ (signe $+$ si erreur, $-$ sinon)\\
\hspace*{1.5em}renormaliser $w$ sur tout le dataset\\[4pt]
\textbf{retourner} $(C_1,\ldots,C_M)$, $(\alpha_1,\ldots,\alpha_M)$
\end{quote}

\textbf{Prédiction.} Pour $x$ et chaque classe $c$ :
\[
  \mathrm{score}_c(x) = \sum_{m=1}^M \alpha_{c,m}\cdot\bigl(2\,C_{c,m}(x)-1\bigr)\cdot\mathbf 1\bigl[(x)_{F_{c,m}}=v_{c,m}\bigr]\cdot\mathbf 1[C_{c,m}\text{ non vide}],
\]
puis $\hat y = \arg\max_c \mathrm{score}_c(x)$ (clause ignorée si son gate n'est
pas satisfait par $x$ ou si elle est vide).



\section*{Validation empirique}

Comparaison de notre modèle à la Tsetlin Machine officielle (\texttt{pyTsetlinMachine},
configurations publiées) et à cinq baselines classiques (\texttt{scikit-learn},
hyperparamètres par défaut). Tous les modèles sont évalués sur les mêmes
fichiers de données et, quand un remélange est nécessaire, la même seed par
run ($20000+\text{run}$). Écart-type rapporté : $\sqrt{\frac1R\sum_r(\mathrm{acc}_r-\bar{\mathrm{acc}})^2}$
(écart-type empirique de population, $1/R$, pas l'estimateur non biaisé
$1/(R{-}1)$). \textbf{Budgets non égalisés} entre notre modèle et la TM
(nombre de clauses différent par dataset, cf.\ ci-dessous) : ces comparaisons
ne sont donc pas encore à budget de calcul égal. PCL (AAAI-25, related work
le plus proche) n'est pas inclus systématiquement : dans nos tests, PCL
s'effondre sur XOR bruité avec ses hyperparamètres publiés (toutes les
clauses restent vides) ; il reste pertinent là où il fonctionne (Iris,
$\approx82\%$ observé séparément).

\subsection*{XOR}
Fichiers fixes (\texttt{NoisyXORTrainingData.txt}/\texttt{TestData.txt},
5000/5000), aucun remélange. Notre modèle : $M{=}200$, $K{=}1$, $S{=}1.0$,
$a{=}0.35$, $N{=}50$, $total{=}5000$. TM : 20 clauses, $T{=}15$, $s{=}3.9$,
200 epochs. 20 runs partout.
\begin{center}
\begin{tabular}{lc}
\toprule
Modèle & Précision \\
\midrule
Notre modèle & $\mathbf{100\%\pm0\%}$ \\
Tsetlin Machine (officielle) & $97.40\%\pm2.23\%$ \\
SVM & $98.12\%\pm0\%$ \\
MLP & $76.72\%\pm1.97\%$ \\
Random Forest & $66.52\%\pm0.30\%$ \\
Régression logistique & $48.94\%\pm0\%$ \\
Naive Bayes & $49.12\%\pm0\%$ \\
\bottomrule
\end{tabular}
\end{center}

\subsection*{Iris}
Fichier unique (\texttt{BinaryIrisData.txt}, 149 exemples), remélangé 80/20 à
chaque run. Notre modèle : $M{=}1500$, $K{=}2$, $S{=}1.0$, $a{=}0.3$,
$N{=}100$, $total{=}20000$. TM : 300 clauses, $T{=}10$, $s{=}3.0$, 500 epochs.
20 runs partout.
\begin{center}
\begin{tabular}{lc}
\toprule
Modèle & Précision \\
\midrule
Notre modèle & $\mathbf{95.5\%\pm4.25\%}$ \\
Random Forest & $95.33\%\pm3.86\%$ \\
MLP & $94.50\%\pm3.38\%$ \\
Tsetlin Machine (officielle) & $93.50\%\pm4.28\%$ \\
SVM & $92.67\%\pm4.03\%$ \\
Régression logistique & $92.50\%\pm4.70\%$ \\
Naive Bayes & $90.67\%\pm4.29\%$ \\
\bottomrule
\end{tabular}
\end{center}
Notre modèle et Random Forest sont dans la marge de bruit l'un de l'autre.

\subsection*{Digits}
Fichiers fixes (\texttt{BinaryDigitsTrainingData192.txt}/\texttt{TestData192.txt},
1437/360), aucun remélange. Notre modèle : $M{=}2000$, $K{=}2$, $S{=}0.05$,
$a{=}1.0$, $N{=}100$, $total{=}20000$. TM : 1000 clauses, $T{=}10$, $s{=}3.0$,
300 epochs. 10 runs partout.
\begin{center}
\begin{tabular}{lc}
\toprule
Modèle & Précision \\
\midrule
Notre modèle & $\mathbf{97.19\%\pm0.40\%}$ \\
SVM & $97.50\%\pm0\%$ \\
Random Forest & $97.31\%\pm0.31\%$ \\
MLP & $97.17\%\pm0.41\%$ \\
Régression logistique & $96.11\%\pm0\%$ \\
Tsetlin Machine (officielle) & $93.69\%\pm0.76\%$ \\
Naive Bayes & $91.39\%\pm0\%$ \\
\bottomrule
\end{tabular}
\end{center}

\subsection*{MNIST}
Fichiers fixes (\texttt{MNISTTraining.txt}/\texttt{Test.txt}, 60000/10000),
aucun remélange. Notre modèle : $M{=}16000$, $K{=}2$, $S{=}1.0$, $a{=}0.3$,
$N{=}100$, $total{=}20000$. TM : 4000 clauses/classe, $T{=}50$, $s{=}10.0$,
20 epochs.
\begin{center}
\begin{tabular}{lc}
\toprule
Modèle & Précision \\
\midrule
SVM & $\mathbf{97.71\%\pm0\%}$ \\
Random Forest & $97.07\%\pm0.06\%$ \\
MLP & $96.90\%\pm1.03\%$ \\
Tsetlin Machine (officielle, 4000 clauses/classe) & $96.62\%\pm0.07\%$ (3 runs) \\
Notre modèle ($M{=}16000$/classe) & $95.85\%$ (1 run) \\
Régression logistique & $91.89\%\pm0\%$ \\
Notre modèle ($M{=}4000$/classe, budget égal TM) & $93.56\%$ (1 run) \\
Naive Bayes & $84.23\%\pm0\%$ \\
\bottomrule
\end{tabular}
\end{center}


\end{document}
